\section{Cumulative span to word span}\label{sec:cumspan_to_wordspan}

Cumulative spanning ($T_N(A) = \MN{D}$, the union of word spans
of all lengths~$\le N$) is weaker than normality
($S_L(A) = \MN{D}$ for a single length~$L$).
An aperiodicity hypothesis upgrades cumulative spanning to a single-length word
span.

\begin{remark}[Aperiodicity condition]\label{rem:aperiodicity}\leanok
    We say $A$ satisfies the \emph{aperiodicity condition}
    if $\Id \in S_1(A)$, i.e.\ the identity matrix lies in the
    span of the Kraus operators $\{A^i\}_{i=0}^{d-1}$.
    This is an additional hypothesis. The one-sided TP normalization
    $\sum_i (A^i)^\dagger A^i = \Id$ does \emph{not} imply
    $\Id \in S_1(A)$; it only gives a quadratic identity among the
    Kraus operators.
    The aperiodicity condition is kept explicit in this section for
    generality.
    For the primitive condition with positive-definite fixed point
    (peripheral-spectrum primitivity together with a positive-definite
    fixed point), normality follows directly from Wolf's primitivity
    equivalence (Theorem~\ref{thm:wolf_6_8_conjunction}) without passing
    through aperiodicity.
\end{remark}

\begin{remark}[Counterexample: cumulative spanning does not imply normality]%
    \label{rem:cumspan_counterexample}\leanok
    In $\MN{2}$, let $A^0 = E_{12}$ and $A^1 = E_{21}$
    (the off-diagonal matrix units).
    Then $\algspn(A) = \MN{2}$ and $T_2(A) = \MN{2}$,
    but the word spans alternate:
    $S_n(A)$ contains only off-diagonal matrices for odd~$n$
    and only diagonal matrices for even~$n$.
    Hence $S_n(A) \neq \MN{2}$ for every~$n$,
    and $A$ is not normal.
    The aperiodicity condition $\Id \in S_1(A)$ fails:
    $S_1(A) = \spn_{\C}\{E_{12}, E_{21}\}$.
\end{remark}

\begin{theorem}[Word span monotonicity from aperiodicity]%
    \label{thm:wordspan_mono_aperiodic}
    \lean{MPSTensor.wordSpan_mono_succ_of_one_mem_wordSpan_one}
    \leanok
    \uses{def:word_span}
    If $\Id \in S_1(A)$, then
    $S_n(A) \le S_{n+1}(A)$ for all~$n$.
\end{theorem}

\begin{proof}\leanok\uses{lem:word_span_mul}
    For any $M \in S_n(A)$,
    $M = M \cdot \Id \in S_n(A) \cdot S_1(A) \subseteq S_{n+1}(A)$
    by Lemma~\ref{lem:word_span_mul}.
\end{proof}

\begin{theorem}[Homogeneous word-span propagation under normalization]%
    \label{thm:wordspan_propagation_unital}
    \lean{MPSTensor.wordSpan_succ_eq_top_of_unital_of_wordSpan_eq_top}
    \lean{MPSTensor.wordSpan_eq_top_of_ge_of_unital}
    \lean{MPSTensor.isNBlkInjective_of_ge_of_unital}
    \leanok
    \uses{def:word_span, def:l_blk_injective}
    Suppose
    \[
        \sum_a A_aA_a^\dagger=\Id
        \qquad\text{and}\qquad
        S_L(A)=\MN{D}.
    \]
    Then, for every \(m\ge L\),
    \[
        S_m(A)=\MN{D}.
    \]
    Equivalently, \(L\)-block injectivity propagates to every larger
    homogeneous length.
\end{theorem}

\begin{proof}\leanok\uses{lem:word_span_mul}
    If \(S_n(A)=\MN{D}\), then for any \(X\in\MN{D}\),
    \[
        X
        =
        \left(\sum_a A_aA_a^\dagger\right)X
        =
        \sum_a A_a(A_a^\dagger X).
    \]
    Since \(A_a^\dagger X\in S_n(A)\), each summand lies in
    \(S_1(A)S_n(A)\subseteq S_{n+1}(A)\). Thus
    \(S_{n+1}(A)=\MN{D}\). Iterating this one-step implication gives the
    result for all \(m\ge L\).
\end{proof}

\begin{theorem}[Cumulative span equals word span under aperiodicity]%
    \label{thm:cumspan_eq_wordspan}
    \lean{MPSTensor.cumulativeSpan_eq_wordSpan_of_one_mem_wordSpan_one}
    \leanok
    \uses{def:cumulative_span, def:word_span}
    If $\Id \in S_1(A)$, then
    $T_n(A) = S_n(A)$ for all~$n$.
\end{theorem}

\begin{proof}\leanok\uses{thm:wordspan_mono_aperiodic}
    By Theorem~\ref{thm:wordspan_mono_aperiodic},
    the word spans are monotone: $S_0 \le S_1 \le \cdots$.
    The supremum $T_n = S_0 + \cdots + S_n$ therefore equals the
    largest term~$S_n$.
\end{proof}

\begin{theorem}[Cumulative spanning plus aperiodicity implies normality]%
    \label{thm:cumspan_to_normal}
    \lean{MPSTensor.isNormal_of_cumulativeSpan_eq_top_of_aperiodic}
    \leanok
    \uses{def:normal, def:cumulative_span, def:word_span}
    If $T_N(A) = \MN{D}$ and $\Id \in S_1(A)$, then $A$ is normal.
\end{theorem}

\begin{proof}\leanok\uses{thm:cumspan_eq_wordspan}
    By Theorem~\ref{thm:cumspan_eq_wordspan},
    $S_N(A) = T_N(A) = \MN{D}$.
\end{proof}

\begin{theorem}[Algebra span plus aperiodicity implies normality]%
    \label{thm:algspan_to_normal}
    \lean{MPSTensor.isNormal_of_algSpan_eq_top_of_aperiodic}
    \leanok
    \uses{def:normal, def:alg_span, def:word_span}
    If $\algspn(A) = \MN{D}$ and $\Id \in S_1(A)$,
    then $A$ is normal.
\end{theorem}

\begin{proof}\leanok\uses{thm:cumspan_to_normal}
    Since $\algspn(A) = \MN{D}$,
    there exists $N$ with $T_N(A) = \MN{D}$
    (by the ascending chain condition on finite-dimensional subspaces).
    Apply Theorem~\ref{thm:cumspan_to_normal}.
\end{proof}

\begin{remark}[Normal canonical form avoids the aperiodicity step]\label{rem:d1_bypasses_aperiodicity}
    The aperiodicity results in Section~\ref{sec:cumspan_to_wordspan} are needed only for the
    block-injective argument through Definition~\ref{def:normal}, where one upgrades cumulative
    spanning to eventual block injectivity. The normal-canonical theory used later in
    Chapters~\ref{ch:canonical}--\ref{ch:bnt} does not pass through this step:
    it works directly with irreducible TP-normalized blocks whose blocked transfer maps
    are primitive and, once the corresponding weighted block data are available,
    places them in normal canonical form in the sense of~%
    \cite{Cirac2016MPDO_arXiv}. There are analogous cumulative-span lemmas for
    blocked irreducible tensors and eigenvector-spreading arguments, but they are
    not needed for the implication stated below.
\end{remark}

\section{Block injectivity from TP/primitive/irreducible blocks}%
\label{sec:ch07_block_inj_from_tp_primitive_irred}

\begin{theorem}[TP-primitive irreducible tensor is injective after blocking]%
    \label{thm:tp_primitive_irred_block_injective}
    \lean{MPSTensor.exists_pos_blockTensor_isInjective_of_tp_primitive_irreducible}
    \leanok
    \uses{thm:normal_tp_primitive_irred,
        lem:one_block_injective_of_injective,
        cor:blk_inj_zero_dim_one}
    Let $A$ be a left-canonical MPS tensor with $D > 0$, irreducible tensor,
    and primitive transfer map.
    Then some positive blocking $A^{[L]}$ is one-site injective.
\end{theorem}

\begin{proof}\leanok
    \uses{thm:normal_tp_primitive_irred,
        lem:one_block_injective_of_injective,
        cor:blk_inj_zero_dim_one}
    If $D=1$, trace preservation gives a nonzero one-site Kraus matrix, and
    every nonzero scalar matrix spans $\MN{1}$; hence $A$ is one-site
    injective, so Lemma~\ref{lem:one_block_injective_of_injective} gives the
    positive witness $L=1$. Suppose now that $D\geq 2$. The normality theorem
    (Theorem~\ref{thm:normal_tp_primitive_irred}) gives a block-injectivity
    witness $L$. If $L=0$, Corollary~\ref{cor:blk_inj_zero_dim_one} would
    force $D=1$, a contradiction. Thus $L>0$, and the blocked-chain
    equivalence identifies this $L$-block injectivity with one-site injectivity
    of the blocked tensor $A^{[L]}$.
\end{proof}

\begin{remark}
    Blocking preserves normality by Theorem~\ref{thm:isNormal_blockTensor}.
\end{remark}
